% Kapitel 4 -- Von Kuonrats Ankunft auf Burg Hardegg und seiner ersten Unterweisung

\section{Von Kuonrats Ankunft auf Burg Hardegg}

\mylettrine{A}{ls} \gls{Kuonrat von Steinare} endlich die Burg \gls{Hardegg} erblickte, stand die Sonne bereits tief über dem Thayatal. Die Burg selbst lag auf dunklem Felshang, und schon von ferne konnte der Knabe die Fackeln sehen, die vom Tor brannten. Aber nicht der Anblick der Burg raubte ihm den Atem, sondern das Getümmel, das sich in dem Burghof abspielte, als er die Zugbrücke überritt.

Denn der Graf \gls{Konrad von Hardegg} war gerade dabei, einen Handel zu besiegeln, der von solcher Art war, dass ihn die Kirche verflucht hätte, hätte sie es je erfahren. Der Fuhrmann \gls{Gruober}, ein von Natur fahler, dürrer Mann mit Augen wie Kohlenglut, war gekommen mit einer Gruppe armseliger Wesen, lauter Weiber, die er in den böhmischen Landen eingefangen hatte, wie der Jäger Hirsche treibt. Sie waren gekleidet in fadenscheinige Stoffe, manche trugen Ketten um die Fußgelenke; ihr Blick war trübe, als hätte ihnen der Menschenwille bereits aus den Augen geleert. So heißt es in den Berichten, die mir erreichbar waren.

Der Graf stand mitten unter ihnen wie ein Geier zwischen Küken, und studierte sie mit Blicken, die dem jungen \gls{Kuonrat von Steinare} böse erscheinen mussten, doch Böses war nicht, was der Knabe in jener Stunde lernen würde zu sehen.

Der Graf winkte den Knappen, näher zu treten. \enquote{Sieh her}, sprach er, \enquote{hier zeigt sich dir die wahre Ordnung der Welt. Ein Weib muss schön sein, verstehst du?} \gls{Kuonrat von Steinare} stand reglos, verstehend und nichtverstehend zugleich. Der Graf deutete auf eine der Gefangenen, eine sehr kleine, sehr dünne Gestalt mit Augen, die noch nicht ganz erloschen waren. \enquote{Sieh diese hier. Schmal wie ein Besen, die Hüften eng, der Leib agil, das ist Schönheit. Nicht aufgedunsen, nicht breit wie ein Mann. Das ist, wofür sich ein Herr müht.}

Dann geschah das, was dem Knaben wie ein Wunder vorkam, oder wie eine Hölle, je nachdem, wie sehr er danach das Gute vom Bösen unterscheiden lernte. Der Graf befahl der Frau, sich zu bücken. Sie erhob den Kopf, und in ihrem Blick flimmerte noch ein Hauch von Widerwille; sie sprach nicht, aber ihr ganzer Körper sagte Nein. Da trat der Graf näher, neigte sich zu ihr hinab, so nah, dass kein andrer hörte, was er flüsterte. Aber der junge \gls{Kuonrat von Steinare} sah, wie das Blut aus ihrem Gesicht wich, wie ihre Lippen sich öffneten stumm in einem Schrei ohne Laut.

Was war in jene Worte eingegangen, dass sie solche Kraft erhielten?

Der Graf richtete sich auf und befahl nochmals: \enquote{Bücke dich.} Und dieses Mal gehorchte die Frau. Ihr Leib beugte sich, ihr Haupt neigte sich, und ihre Hände falteten sich hinter dem Rücken, als wäre die Seele aus ihr gewichen und durch Gehorsam ersetzt worden. Der Graf sah sich um, prüfend, als wolle er sichergehen, dass der Knabe dies alles sehe. Und dann, dies ist das wahrste und das schrecklichste, das ich zu berichten habe, ergriff er die \gls{Geisslein}, die an seiner Hüfte hing, ein \gls{Geisslein} aus Leder, rau und verflochten, und schlug die Frau damit genau zehnmal auf den Hintern, während sie unbeweglich blieb.

\enquote{So}, sprach der Graf mit kalter Regelmäßigkeit, der jedes Mal eine Art Wohlklang unterlag, als rechne er auf seinen Fingern, \enquote{das ist die Buße für ihre Hartnäckigkeit. Wer nicht von selbst gehorcht, muss den Weg durch Schmerz lernen.} Die Schläge waren nicht so heftig, dass die Frau zu Blut aufgerissen wäre, doch war es kein Spiel. Es war Ausübung von Macht, und Macht, die keines Wortes mehr bedurfte.

Als es vorbei war, sackte die Frau zusammen. Der Graf aber wandte sich dem Knaben zu und lächelte, ein Lächeln ohne Wärme. \enquote{Siehst du, \gls{Kuonrat von Steinare}? Mit Worten allein, mit den rechten Worten, kann man ein Weib brechen, ehe man auch nur einen Finger rührt. Ich habe ihr gesagt, was ihr widerfährt, wenn sie nicht gehorcht, und ihr Leib hat verstanden, bevor der Verstand noch etwas vermochte.} Der Knabe aber saß auf seinem Pferde, stumm und ergriffen, und dachte, der Graf habe die Frau nur mit Befehlen bezwungen, mit nichts weiter als mit Stimme und Wort. So wissen die Unerfahrenen nichts von der Tiefe der Bosheit.
